\documentclass{article}
\usepackage{neurips_2026}
\usepackage[utf8]{inputenc}
\usepackage[T1]{fontenc}
\usepackage{hyperref}
\usepackage{url}
\usepackage{booktabs}
\usepackage{amsfonts,amsmath,amsthm,amssymb}
\usepackage{nicefrac}
\usepackage{microtype}
\usepackage{xcolor}
\usepackage{graphicx}
\usepackage{multirow}
\newtheorem{proposition}{Proposition}
\newtheorem{corollary}{Corollary}
\theoremstyle{remark}\newtheorem{remark}{Remark}

\title{How Much Room Is Left Below PPI++? Efficiency Limits and Cross-Fitted Recalibration for Prediction-Powered Mean Inference When Labels, Unlabeled Data and Predictor Are Fixed}
\author{Anonymous Author(s)}

\begin{document}
\maketitle

\begin{abstract}
Prediction-powered inference (PPI) builds valid confidence intervals from a small labeled sample, a large unlabeled sample and a fixed pre-trained predictor. We ask a narrow question: with all three held fixed and nominal coverage required, how much shorter than PPI/PPI++ can an interval for the population mean be, and when? We show that, if the predictor output is the only information used, the answer is exactly the non-affine part of the calibration curve $\mathbb E[Y\mid f(X)]$: PPI++ is asymptotically efficient if and only if this curve is affine (or there is no unlabeled data), and the width reduction available to any estimator is at most $1-\sqrt{1-sN/(n+N)}$, where $s$ is the share of residual variance that is non-affine. A cross-fitted recalibrated estimator (isotonic, spline or CV-selected) attains the bound and stays valid for any calibration fit, without convergence rates. In Monte-Carlo experiments with known truth (2000 replicates per synthetic cell, 1000 per real-data cell), on seven example data sets from \texttt{ppi-python} the attainable gain over PPI++ is at most 1.2\% in width (nonlinear share $\le 2.5\%$) and the recalibrated intervals realise essentially none of it: at $n=500$ they are 0.5--2\% wider than PPI++ on six sets and 0.6\% narrower on one, and up to 19\% wider at $n=50$. Gains appear when the predictor outputs are not on the scale of the label (logit or otherwise monotonically distorted scores): 2--27\% shorter on distorted versions of the same data, within about 2 percentage points of the bound. A hard binary predictor offers no headroom, and under distribution shift the recalibrated intervals lose coverage completely where PPI keeps it. Separately, ppi\_py's clipping of the PPI++ weight costs about 5\% width on one data set. Our conclusion is conditional: PPI++ is close to optimal for well-calibrated predictors, and recalibration is worth it only for predictors whose scores are visibly non-affine in the label.
\end{abstract}

\section{Introduction}
In ecology, astronomy and social science a typical analysis has three ingredients: a small set of $n$ records with expensive gold-standard labels, a much larger set of $N$ records with features only, and a pre-trained model that predicts labels for any record. Prediction-powered inference (PPI) \citep{angelopoulos2023ppi} and its power-tuned version PPI++ \citep{angelopoulos2023ppipp} use the predictions to shrink confidence intervals while keeping nominal coverage, regardless of how good the predictor is. Subsequent work changes the inputs: it trains the predictor on the labels with cross-fitting \citep{zrnic2023cross}, spends the label budget more cleverly \citep{zrnic2024active}, or stratifies the labeled set \citep{fisch2024stratified}. Here we fix everything: the labeled set, the unlabeled set and the predictor are given, and nothing can be added.

\textbf{Question.} \emph{Given exactly these inputs, can the interval be made shorter than PPI/PPI++ without losing coverage; by how much; and under which conditions?}

We study the population mean $\theta=\mathbb E[Y]$. It is the simplest target for which a sharp answer is available and a closed-form efficiency analysis is possible; it is also the target of every example in the \texttt{ppi-python} data used here. Our answer has three parts.
\begin{enumerate}
\item \textbf{A tight characterization (Sec.~\ref{sec:theory}).} If the predictor output $Z=f(X)$ is the only information used, the asymptotic variance of any regular estimator is at least $\mathbb E\mathrm{Var}(Y|Z)/n+\mathrm{Var}(\mathbb E[Y|Z])/(n+N)$. PPI++ attains it iff $\mathbb E[Y|Z]$ is affine in $Z$ or $N/(n+N)\to0$; otherwise the variance gap is $\mathrm{Var}(h)\,N/(n(n+N))$ with $h$ the non-affine part of the calibration curve. A hard 0/1 predictor can never be improved. This is the semiparametric theory of missing-at-random means \citep{robins1994estimation,tsiatis2006semiparametric} applied to PPI; it is also close to what is known for semi-supervised mean estimation \citep{zhang2019semi}, so we do not claim novelty for the bound itself but for its use as a yardstick for PPI in the fixed-inputs regime.
\item \textbf{A route to the bound (Sec.~\ref{sec:method}).} A cross-fitted recalibrated PPI interval whose validity needs only $L^2$-stability of the calibration fit, not a rate, and that is therefore valid for \emph{any} calibration learner (isotonic, spline, bins, CV-selected).
\item \textbf{Evidence about when it pays (Sec.~\ref{sec:exp}).} Across synthetic data with known nonlinearity and seven real data sets resampled with known truth, we measure coverage and width with uncertainty, and compare with the closed-form bound. The headline is negative for the common case and positive for a specific, identifiable case: gains require a predictor whose scores are non-affine in the label; on the unmodified predictors that ship with \texttt{ppi-python} there is essentially nothing left to gain from the predictor output.
\end{enumerate}

\section{Related work}
\textbf{Prediction-powered inference.} PPI \citep{angelopoulos2023ppi} gives valid intervals for means, quantiles and regression coefficients from a labeled set, an unlabeled set and a black-box predictor; PPI++ \citep{angelopoulos2023ppipp} adds a data-driven weight $\lambda$ that guarantees (asymptotically) no loss relative to the classical interval. Cross-PPI \citep{zrnic2023cross} handles predictors trained on the same labels; active inference \citep{zrnic2024active} chooses which points to label; stratified PPI \citep{fisch2024stratified} allocates labels across strata of the predictor. All of these change a part of the input or the sampling design. Our setting excludes that; the closest ingredient is stratification on the predictor, which we include as a post-stratification baseline (bin10). \citet{miao2023popinf} propose assumption-lean, data-adaptive post-prediction inference, another way of using ML predictions for valid inference with a small labeled set.

\textbf{Semi-supervised and missing-data inference.} Estimating a mean from a labeled and an unlabeled sample has a classical efficiency theory \citep{robins1994estimation,bickel1993efficient,tsiatis2006semiparametric}, developed for the semi-supervised mean in \citet{zhang2019semi}, with high-dimensional covariates in \citet{zhang2022high} and M-estimation in \citet{song2024general}. PPI++ for the mean is a member of the class of estimators that are linear in the predictor; what has not, to our knowledge, been laid out is how large the gap to the bound is for the predictors and data sets on which PPI is demonstrated, or its closed form in PPI's notation. Survey statistics has long used nonparametric model-calibrated estimators with the same structure \citep{sarndal1992model,montanari2005nonparametric}, and double machine learning provides the cross-fitting device \citep{chernozhukov2018dml}.

\textbf{Calibration of classifiers.} That predicted scores need not be on the scale of the label is well known: boosted trees and similar models produce sigmoid-shaped distortions of probabilities \citep{niculescu2005predicting} and modern networks are over-confident \citep{guo2017calibration}. This is exactly the situation in which a linear weight $\lambda$ cannot correct the predictor.

\section{Setting and what is attainable when everything is fixed}\label{sec:theory}
\paragraph{Setting.}
Let $(X_i,Y_i)$, $i=1,\dots,m$, be i.i.d.\ with $m=n+N$. A uniformly random subset $L$ of size $n$ is labeled; for $i\notin L$ the label is missing (missing completely at random, MCAR, by design of the labeling). The pre-trained predictor $f$ is fixed and not trained on any of these data; write $Z_i=f(X_i)\in\mathbb R$ (the argument below only uses that $Z_i$ is a measurable function of $X_i$). The analyst sees $Z_i$ for all $m$ points and $Y_i$ for $i\in L$, and wants a $(1-\alpha)$ interval for $\theta=\mathbb E[Y]$. We chose the mean because it is the simplest target for which the efficiency question has a closed-form answer, it is the target for which PPI was first demonstrated, and it is what all nine example data sets in \texttt{ppi-python} are used for. Write
\[
g(z)=\mathbb E[Y\mid Z=z],\quad \sigma^2=\mathbb E\,\mathrm{Var}(Y\mid Z),\quad V=\mathrm{Var}(Z),\quad \beta=\tfrac{\mathrm{Cov}(Y,Z)}{V},\quad h(z)=g(z)-\mathbb E g-\beta (z-\mathbb E Z),
\]
so $h$ is the part of the calibration curve $g$ that is \emph{not} affine in the predictor output, and $v_h=\mathrm{Var}\,h(Z)$. Let $s=v_h/(\sigma^2+v_h)$ be the \emph{nonlinear share} of the residual variance left by the best affine recalibration. All variances below are asymptotic variances of the estimator of $\theta$ (i.e.\ squared interval half-widths up to the factor $z_{1-\alpha/2}^2$); $n,N\to\infty$ with $n/m\to\pi\in(0,1]$.

\paragraph{Existing estimators.}
PPI with weight $\lambda$ uses $\hat\theta_\lambda=\frac1N\sum_{j\notin L}\lambda Z_j+\frac1n\sum_{i\in L}(Y_i-\lambda Z_i)$ \citep{angelopoulos2023ppi,angelopoulos2023ppipp}; $\lambda=1$ is the original PPI, $\lambda=0$ the classical sample mean, and PPI++ picks $\lambda$ from the data to minimize variance. Its asymptotic variance is $\mathrm{Var}(Y-\lambda Z)/n+\lambda^2V/N$.

\begin{proposition}[Variance of the PPI family]\label{prop:ppi}
Minimizing over $\lambda\in\mathbb R$ gives $\lambda^\star=\beta N/m$ and
$V_{\mathrm{PPI++}}=\dfrac{\sigma^2+v_h}{n}+\dfrac{\beta^2V}{m}$. For fixed $\lambda$,
$V_\lambda-V_{\mathrm{PPI++}}=V\frac{m}{nN}(\lambda-\lambda^\star)^2$; in particular PPI with $\lambda=1$ is suboptimal by $V\frac{m}{nN}(1-\beta N/m)^2$, which vanishes only if $\beta N/m=1$.
\end{proposition}

\begin{proposition}[Efficiency bound with the predictor output as the only information]\label{prop:bound}
In the model where $(Z,Y)$ has an unrestricted distribution (finite variance), any regular asymptotically linear estimator of $\theta$ based on $\{Z_i\}_{i\le m}$ and $\{Y_i\}_{i\in L}$ has asymptotic variance at least
\[V_{\mathrm{eff}}=\frac{\sigma^2}{n}+\frac{\mathrm{Var}\,g(Z)}{m}=\frac{\sigma^2}{n}+\frac{\beta^2V+v_h}{m}.\]
It is attained by $\hat\theta=\frac1n\sum_{L}(Y_i-g(Z_i))+\frac1m\sum_{i\le m}g(Z_i)$, whose influence function is $\frac{R}{\pi}(Y-g(Z))+g(Z)-\theta$.
\end{proposition}
\begin{proof}[Proof sketch]
This is the standard semiparametric result for the mean under MCAR missingness with known labeling probability $\pi$ \citep{robins1994estimation,bickel1993efficient,tsiatis2006semiparametric}: with the nonparametric model for $(Z,Y)$, the observed-data influence function of $\theta$ is unique and equals $\frac{R}{\pi}(Y-\theta)-\frac{R-\pi}{\pi}\mathbb E[Y-\theta\mid Z]$. Its variance per observation is $\mathbb E(Y-g)^2/\pi+\mathrm{Var}\,g$, which in units of the sample sizes is the displayed expression. The case $Z=\text{const}$ recovers $\mathrm{Var}(Y)/n$, and $N=\infty$ recovers $\sigma^2/n$ \citep{zhang2019semi}.
\end{proof}

\begin{corollary}[When PPI++ cannot be improved]\label{cor:gap}
$V_{\mathrm{PPI++}}-V_{\mathrm{eff}}=v_h\,\dfrac{N}{n\,m}$, hence
\[
\frac{V_{\mathrm{eff}}}{V_{\mathrm{PPI++}}}=1-\frac{N}{m}\cdot\frac{s}{1+\kappa}\ \ \ge\ 1-\frac{N}{m}\,s,\qquad \kappa=\frac{n}{m}\frac{\beta^2V}{\sigma^2+v_h}\ge 0 .
\]
Thus (i) PPI++ is asymptotically efficient if and only if $g$ is a.s.\ affine in $Z$ ($v_h=0$) or $N/m\to0$ (essentially no unlabeled data); (ii) whenever the predictor output takes two values (a hard 0/1 prediction), $g$ is automatically affine, so no estimator that uses only $Z$ can beat PPI++; (iii) the interval-width reduction is at most $1-\sqrt{1-sN/m}\le 1-\sqrt{1-s}$, e.g.\ about $5\%$ for $s=0.10$ and about $1\%$ for $s=0.02$.
\end{corollary}
The corollary follows by subtracting the two expressions in Props.~\ref{prop:ppi} and~\ref{prop:bound}. It says exactly \emph{what} room remains below PPI++ with labeled set, unlabeled set and predictor all fixed: only the non-affine part of the calibration curve $\mathbb E[Y\mid f]$, discounted by the fraction $N/m$ of unlabeled data. It also shows that the large-$N$ regime that motivates PPI is where the headroom is largest. Using side features $X$ beyond $Z=f(X)$ lowers the bound further to $\mathbb E\,\mathrm{Var}(Y|X)/n+\mathrm{Var}(\mathbb E[Y|X])/m\le V_{\mathrm{eff}}$, but attaining it requires learning $\mathbb E[Y|X]$ from the labels, i.e.\ in effect replacing the predictor; we only quantify how large this extra gap is on the data at hand (\S\ref{sec:exp-x}).

\begin{remark}[Invariance]
A monotone reparametrization $Z\mapsto T(Z)$ (e.g.\ logits instead of probabilities) leaves $g$-based estimators with monotone or nonparametric $g$ unchanged, but changes $\beta$, $v_h$ and hence the width of PPI/PPI++. A predictor whose scores are not on the scale of $Y$ is therefore exactly the case with $v_h>0$.
\end{remark}

\section{Cross-fitted recalibrated PPI}\label{sec:method}
Proposition~\ref{prop:bound} suggests the obvious route to shorter intervals: estimate the calibration curve $g=\mathbb E[Y|Z]$ from the labeled data instead of fixing it to $z\mapsto\lambda z$. The estimator is an augmented-inverse-probability-weighted estimator with known labeling probability \citep{robins1994estimation}, equivalently a model-assisted/model-calibrated survey estimator with a nonparametric working model \citep{sarndal1992model,montanari2005nonparametric}; it is also the mean-estimation special case of cross-fitted DML \citep{chernozhukov2018dml}. We do not claim the estimator is new; the point of this paper is to measure how much it can gain over PPI/PPI++ and when.

\paragraph{Algorithm.} Split the labeled indices into $K=5$ random folds $L_1,\dots,L_K$. For each $k$ fit $\hat g_{-k}$ on $L\setminus L_k$ by a one-dimensional regression of $Y$ on $Z$. Put $\hat g_i=\hat g_{-k(i)}(Z_i)$ for $i\in L$ (out-of-fold) and $\bar g_j=\frac1K\sum_k\hat g_{-k}(Z_j)$ for unlabeled $j$. Then
\[
\hat\theta=\frac1n\sum_{i\in L}(Y_i-\hat g_i)+\frac1m\Big(\sum_{i\in L}\hat g_i+\sum_{j\notin L}\bar g_j\Big),\qquad
\hat V=\frac{\widehat{\mathrm{Var}}_L(Y-\hat g)}{n}+\frac{\widehat{\mathrm{Var}}_{\text{all}}(\hat g)}{m},
\]
and the interval is $\hat\theta\pm z_{1-\alpha/2}\sqrt{\hat V}$ (sample variances with $n{-}1$ and $m{-}1$ denominators). With $\hat g\equiv\lambda z$ this is PPI up to the use of the labeled predictions in the second mean; with $\hat g$ the OLS line it is the survey regression estimator, asymptotically equivalent to power-tuned PPI. Learners (all in \texttt{code/cppi.py}): \textbf{lin} (OLS line); \textbf{iso} (isotonic regression, increasing, \citealp{barlow1972isotonic}); \textbf{bin10} (10 equal-mass bins, i.e.\ post-stratification on the predictor output, in the spirit of stratified PPI \citealp{fisch2024stratified} but with random rather than allocated labels); \textbf{spl} (natural cubic regression spline with 5 knots at the 5\%--95\% quantiles of the training scores, \citealp{hastie2009elements}); \textbf{auto} (choose among lin, iso, spl by inner 4-fold CV mean squared error, nested inside each outer fold, so that selection never sees the points it is evaluated on). No tuning on the evaluation data was done; $K$, knot counts and bin counts were fixed before the experiments.

\begin{proposition}[Validity]\label{prop:valid}
Suppose $\mathbb E Y^2<\infty$ and there is a function $g^*$ with $\mathbb E g^*(Z)^2<\infty$ such that $\mathbb E[(\hat g_{-k}(Z)-g^*(Z))^2\mid\text{training fold}]\to_p0$ for each $k$, where $Z$ is an independent copy. Then $(\hat\theta-\theta)/\sqrt{\hat V}\Rightarrow N(0,1)$, so the interval has asymptotic coverage $1-\alpha$. Moreover $\hat\theta$ is exactly conditionally unbiased given the training folds for \emph{any} $\hat g$. If $g^*=g$ the interval attains $V_{\mathrm{eff}}$.
\end{proposition}
The limit $g^*$ need not be the truth: a misspecified learner (e.g.\ a monotone fit to a non-monotone $g$) yields a valid but longer interval, with variance $\mathbb E(Y-g^*)^2/n+\mathrm{Var}\,g^*(Z)/m$. This is the same robustness-to-the-predictor that makes PPI valid for any $f$; the proof (Appendix~\ref{app:proof}) uses only conditional Chebyshev bounds and no convergence \emph{rate} for $\hat g$. The guarantee is asymptotic; Sec.~\ref{sec:exp} measures finite-sample coverage down to $n=25$. Data-driven learner selection (\textbf{auto}) is covered by the proposition when the selected learner's limit is stable; we do not prove this and rely on simulation.

\section{Experiments}\label{sec:exp}
\paragraph{Protocol.} Every number comes from Monte-Carlo replicates with a known truth. \emph{Real-data cells:} the seven \texttt{ppi-python} example sets that contain $(Y,\hat Y)$ for all records (Galaxies $P{=}16{,}743$, Forest $1{,}596$, AlphaFold $10{,}802$, Gene expression $61{,}150$, Census income $380{,}091$, Census health insurance $318{,}215$, Census education $318{,}215$ records; $P$ is the population size) are each treated as the population; a replicate draws $n$ labeled and $N$ unlabeled records i.i.d.\ \emph{with replacement} and the truth is the population mean of $Y$. The predictor is the one shipped with the data, never refit. (Records repeat when $n+N$ is not small relative to $P$, most visibly for Forest; the DGP is still i.i.d.\ from the empirical distribution, but the efficiency bound of Prop.~\ref{prop:bound} for that DGP is then not the one we compute from cross-validated fits, so we flag Forest at $n\ge1000$.) \emph{Synthetic cells:} (a) $f\sim N(0,1)$, $Y=f+c\,(f^2-1)/\sqrt2+\varepsilon$, $\varepsilon\sim N(0,1)$, so $s=c^2/(1+c^2)$ is known exactly; (b) a classifier that outputs logits $f\sim N(\mu,\tau^2)$ with $Y\sim\mathrm{Bernoulli}(\mathrm{expit}(f))$ (so $\mathbb E[Y|f]=\mathrm{expit}(f)$ is non-affine), with $(\mu,\tau)=(-1,2)$ ($s=0.06$) and $(-1,4)$ ($s=0.23$). \emph{Methods:} classical; PPI ($\lambda{=}1$) and PPI++ from \texttt{ppi-python} 0.2.3 (which clips $\hat\lambda$ to $[0,1]$); PPI++ with the unclipped plug-in $\hat\lambda=\widehat{\mathrm{Cov}}(Y,f)/((1+n/N)\widehat{\mathrm{Var}}(f))$, which is the best member of the PPI family and is our reference; and cross-fitted lin / bin10 / iso / spl / auto (Sec.~\ref{sec:method}, $K=5$). All intervals are 95\% Wald-type. \emph{Metrics:} coverage with Wilson intervals; width reported as the paired per-replicate ratio to the reference (mean $\pm1.96$ s.e.\ over replicates, so values that differ from 1 by more than the $\pm$ are not noise); the closed-form ratio $\sqrt{V_{\rm eff}/V_{\rm PPI++}}$ from Props.~\ref{prop:ppi}--\ref{prop:bound} (``bound'') computed from the population (real data: 5-fold cross-validated isotonic/spline fit of $\mathbb E[Y|f]$, so it is itself an estimate, and is capped by the linear fit). Replicates: 2000 per synthetic cell, 1000 per real-data cell. Seeds are derived from (experiment cell, replicate index), so the same labeled/unlabeled draws are shared across methods within a cell.

\subsection{Synthetic data: the bound is tight, and recalibration reaches it}
Figure~\ref{fig:theory} varies the nonlinear share $s$ at $n=500$. The cross-fitted linear interval tracks PPI++ (within 1.2\%, the small price of cross-fitting), while spline/auto intervals follow the efficiency bound: for $N{=}50{,}000$ the width ratio to PPI++ is 0.47 at $s=0.8$, 0.72 at $s=0.5$, 0.91 at $s=0.2$ against the bound 0.46, 0.71, 0.90. Over the 20 cells the empirical ratio of spline/auto exceeds the bound by $0.010\pm0.003$ (mean $\pm$ s.d. over cells; maximum 0.015), which is the cost of estimating the calibration curve. Gains vanish when $N$ is small ($N{=}100$: at most 5.5\% at $s=0.8$), as Cor.~\ref{cor:gap} predicts. Isotonic regression falls short here because $c(f^2-1)$ is non-monotone: its limit $g^*\neq g$, so the interval stays valid but longer (0.67 vs 0.47), as Prop.~\ref{prop:valid} allows. Coverage of auto was between 94.2\% and 95.9\% in all 20 cells (the two cells where several methods sit at 93.4--94.0\% are shared by PPI and PPI++ because the replicates are shared; Table~\ref{tab:valid}).

\begin{figure}[t]\centering
\includegraphics[width=\linewidth]{../figs/fig_theory.pdf}
\caption{Width of cross-fitted intervals relative to PPI++ (unclipped $\lambda$) versus nonlinear share $s$ of the calibration curve ($n=500$; synthetic quadratic DGP; 2000 replicates; bars: $\pm1.96$ s.e.). Dashed: closed-form bound $\sqrt{V_{\rm eff}/V_{\rm PPI++}}$.}\label{fig:theory}
\end{figure}

\paragraph{Small samples (Fig.~\ref{fig:smalln}).} The price of flexibility is paid at small $n$. With a linear truth ($c=0$), auto is 14.5\%, 7.1\%, 3.4\%, 0.6\% wider than PPI++ at $n=25,50,100,500$; spline alone is 65\% wider at $n=25$. With a quadratic truth ($c=1$) auto breaks even at about $n=30$ and gains $10.1\%$ ($n=50$), 20\% ($n=100$), 27\% ($n=500$). Wald coverage at $n=25$ is 91--95\% across methods (classical 92.7--93.0\%); the cross-fitted intervals are not worse and, at $n=25$--$50$, better than ppi\_py's PPI++ (e.g.\ $c=1$, $n=50$: 93.2\% auto vs 91.8\% PPI++, 92.0\% for the $t$-version of unclipped PPI++ vs 93.6\% for auto). A Student-$t$ calibration of the half-widths with $n-1$ d.f.\ narrows but does not close this gap for PPI++ (92.0\% at $n=50$).

\begin{figure}[t]\centering
\includegraphics[width=.82\linewidth]{../figs/fig_smalln.pdf}
\caption{Small-$n$ behaviour ($N=20{,}000$, quadratic synthetic DGP, 2000 replicates). Top: width relative to PPI++ (unclipped). Bottom: coverage with Wilson 95\% intervals.}\label{fig:smalln}
\end{figure}

\subsection{Real data with the shipped predictors: essentially nothing to gain}\label{sec:exp-real}
Table~\ref{tab:real} (and Appendix Figs./Tables) shows the central negative result. The nonlinear share of the calibration curve is at most $2.5\%$ (Gene expression) and $\le0.5\%$ for the other six data sets; the implied bound on the achievable width reduction over PPI++ at $n=500$, $N=20{,}000$ is 0.9\% (Gene expression), 0.2\% (Census income) and $\le0.1\%$ elsewhere. The cross-fitted auto interval does not realise even that: at $n=500$ it is 0.5--2.1\% \emph{wider} than PPI++ on six data sets, and 0.6\% narrower on Gene expression; at $n=2000$ all six are within 0.6\% of PPI++ (the 1\% ``gains'' on Forest at $n\ge1000$ are the resampling artifact flagged above); at $n=50$ auto is 3--19\% wider except on Gene expression ($-2.4\%$). At $n\ge500$ the coverage of every method lies between 91.5\% and 96.8\% (mean 94.5--95.2\%); the lowest cells (92.2--92.8\%, Census income with $N=5000$) are shared by all methods including the classical interval. At smaller $n$ the Census-income coverage is 82--92\% for $n\le200$ for every method (heavy-tailed income), so these shortfalls come from the Wald approximation and not from recalibration.

Two other facts matter for practice. (i) The headroom \emph{relative to the classical interval} is large (width 1.2--3.2 times higher), so using a predictor is worthwhile; the headroom beyond PPI++ is not. (ii) The choice \emph{within} the PPI family matters more than recalibration: PPI with $\lambda=1$ is 4.9\% wider than PPI++ on Gene expression ($\hat\beta=1.15$) and 1.8\% on Census income ($\hat\beta=1.38$), and ppi\_py's PPI++ shows the same loss because its $\hat\lambda$ is clipped at 1 while the unclipped optimum is above 1 (bound ratio of the clipped version to the unclipped one: 1.056 and 1.018).

\begin{table}[t]\centering\small
\caption{Real data, shipped predictors, $n=500$, $N=20{,}000$, 1000 replicates. corr: correlation of $Y$ and $\hat Y$; $\hat\beta$ slope of $Y$ on $\hat Y$; $s$: cross-validated nonlinear share of residual variance; widths are ratios to PPI++ with unclipped $\hat\lambda$ ($\pm$: 1.96 s.e.); ``bound'' is $\sqrt{V_{\rm eff}/V_{\rm PPI++}}$ estimated from the population. Forest has only 1596 records ($n/P=0.31$).}\label{tab:real}
\resizebox{\linewidth}{!}{\begin{tabular}{lrrr rrrr rrrr}
\toprule
 & & & & \multicolumn{4}{c}{coverage (\%)} & \multicolumn{4}{c}{width / PPI++ (unclipped)}\\
\cmidrule(lr){5-8}\cmidrule(lr){9-12}
data set & corr & $\hat\beta$ & $s$ & classical & PPI & PPI++ & auto & classical & PPI & auto & bound\\
\midrule
Galaxies & 0.74 & 0.98 & 0.0\% & 95.3 & 95.1 & 94.7 & 95.2 & 1.47 & 1.00 & 1.007$\;{\scriptstyle\pm0.000}$ & 1.000\\
Forest & 0.54 & 0.90 & 0.0\% & 96.1 & 95.3 & 95.1 & 95.1 & 1.18 & 1.00 & 1.005$\;{\scriptstyle\pm0.001}$ & 1.000\\
AlphaFold & 0.82 & 1.03 & 0.2\% & 95.0 & 94.8 & 94.3 & 94.2 & 1.69 & 0.99 & 1.005$\;{\scriptstyle\pm0.001}$ & 0.999\\
Gene expression & 0.96 & 1.15 & 2.5\% & 94.1 & 93.9 & 94.0 & 93.8 & 3.24 & 1.05 & 0.994$\;{\scriptstyle\pm0.001}$ & 0.991\\
Census income & 0.61 & 1.38 & 0.5\% & 94.2 & 93.5 & 93.1 & 93.4 & 1.26 & 1.02 & 1.021$\;{\scriptstyle\pm0.002}$ & 0.998\\
Census health ins. & 0.71 & 0.99 & 0.1\% & 94.2 & 94.7 & 94.8 & 95.0 & 1.40 & 1.00 & 1.008$\;{\scriptstyle\pm0.000}$ & 1.000\\
Census education & 0.63 & 0.98 & 0.0\% & 93.9 & 94.7 & 94.1 & 94.5 & 1.28 & 1.00 & 1.011$\;{\scriptstyle\pm0.001}$ & 1.000\\
\bottomrule
\end{tabular}}
\end{table}

\subsection{Distorted predictor outputs: where recalibration pays}\label{sec:exp-distort}
Classifiers often emit scores that are monotone but not affine in the probability of the label (logits, over-confident or squashed scores \citealp{niculescu2005predicting,guo2017calibration}). To measure what that costs, we apply a fixed monotone map to the shipped predictor outputs (logit of the probability score; cubed score; $\exp(\hat y/4)$ for Gene expression; $\sqrt{\cdot}$ for Census income) and keep everything else. These are \emph{constructed} cases, chosen to represent the miscalibration mechanism, not evidence of how common it is. Table~\ref{tab:distort} shows that the width reduction of auto/iso relative to PPI++ ranges from $-2\%$ to $-27\%$, matching the bound within about 2 percentage points ($n=1000$), with nominal coverage; plain PPI ($\lambda{=}1$) becomes 4.9--9.5 times wider than PPI++ for logit scores (and 4.6 times for the exponential distortion) because the scores are no longer on the scale of $Y$. With logit-scale synthetic scores of moderate sharpness ($s=0.06$) the gain is only about 2\%, with sharper logits ($s=0.23$) it is 10\% at $n=500$ (bound 11.7\%) and 3\%, 8\%, 10\% at $n=50,100,200$ (Appendix Table~\ref{tab:logit}).

\begin{table}[t]\centering\small
\caption{Monotonically distorted predictor outputs, $n=1000$, $N=20{,}000$, 1000 replicates. Width ratios to PPI++ with unclipped $\hat\lambda$. The last column is the ratio of ppi\_py's clipped PPI++ to the unclipped one.}\label{tab:distort}
\resizebox{\linewidth}{!}{\begin{tabular}{l r rrrr rrr r}
\toprule
 & & \multicolumn{4}{c}{width / PPI++ (unclipped)} & \multicolumn{3}{c}{coverage (\%)} & \\
\cmidrule(lr){3-6}\cmidrule(lr){7-9}
predictor output & $s$ & PPI & iso & auto & bound & PPI++ & iso & auto & $\lambda^{\rm clip}$-PPI++\\
\midrule
Galaxies, logit & 9\% & 8.46 & 0.964$\;{\scriptstyle\pm0.001}$ & 0.963$\;{\scriptstyle\pm0.001}$ & 0.957 & 95.1 & 94.3 & 94.2 & 1.00\\
Galaxies, pow3 & 13\% & 1.00 & 0.945$\;{\scriptstyle\pm0.001}$ & 0.942$\;{\scriptstyle\pm0.001}$ & 0.937 & 95.1 & 95.2 & 96.0 & 1.00\\
AlphaFold, logit & 23\% & 9.50 & 0.900$\;{\scriptstyle\pm0.001}$ & 0.899$\;{\scriptstyle\pm0.001}$ & 0.892 & 94.7 & 95.0 & 95.0 & 1.00\\
Forest, logit & 5\% & 4.94 & 0.975$\;{\scriptstyle\pm0.001}$ & 0.977$\;{\scriptstyle\pm0.001}$ & 0.979 & 94.7 & 95.1 & 95.2 & 1.00\\
Gene expression, exp & 60\% & 4.62 & 0.738$\;{\scriptstyle\pm0.001}$ & 0.731$\;{\scriptstyle\pm0.001}$ & 0.731 & 94.8 & 95.4 & 95.3 & 1.00\\
Census health ins., logit & 8\% & 6.50 & 0.971$\;{\scriptstyle\pm0.001}$ & 0.967$\;{\scriptstyle\pm0.001}$ & 0.962 & 94.4 & 94.3 & 94.5 & 1.00\\
Census income, sqrt & 8\% & 1.20 & 0.987$\;{\scriptstyle\pm0.002}$ & 0.976$\;{\scriptstyle\pm0.001}$ & 0.961 & 94.6 & 95.5 & 94.9 & 1.20\\
\bottomrule
\end{tabular}}
\end{table}

\subsection{Binary predictions, distribution shift, and a fixed-sets application}\label{sec:exp-shift}
\paragraph{Binary predictor.} The plankton predictions are 0/1. Then $g$ is affine and Cor.~\ref{cor:gap}(ii) applies: auto/iso/lin are within 0.1\% of PPI++ in width for $n\ge500$ and 1.2\% narrower at $n=100$ (coverage in Table~\ref{tab:shift}, top). The same holds for the ballots application below.
\paragraph{Distribution shift.} The labeled plankton records are from 2013 and the unlabeled ones from 2014; the target is the 2014 mean. PPI's rectifier $Y-\hat Y$ is approximately stable across years, so PPI ($\lambda{=}1$) keeps 94\% coverage until $n=2000$ and degrades at $n=10{,}000$ (91\%), whereas every interval that estimates the calibration curve on 2013 labels (PPI++ with $\hat\lambda<1$, lin, iso, auto) collapses: auto covers 84\% at $n=500$, 9\% at $n=2000$, 0\% at $n=10{,}000$. Efficiency comes from leaning harder on the stationarity of $\mathbb E[Y|f]$, which is exactly what fails under shift. Estimating the calibration curve thus makes the interval \emph{less} robust than the PPI default; we have no remedy and make no claim about shift-robust variants.
\paragraph{Fixed-sets application.} For the ballots data ($n=1027$, $N=78{,}150$, binary predictions, no ground truth for the unlabeled records) the classical interval has width 0.0591, PPI 0.0133, PPI++ 0.0131 and all cross-fitted intervals 0.0132 (Appendix Table~\ref{tab:ballots}); with a binary predictor nothing below PPI++ is possible, in line with the theory. Here $\hat Y=1\Rightarrow Y=1$ for all labeled records, so the predictor is nearly perfect on one side.

\begin{table}[t]\centering\small
\caption{Coverage (\%) on plankton (1000 replicates, 95\% nominal). Top: labeled and unlabeled data from the same year (2013); bottom: labeled 2013, unlabeled 2014, target is the 2014 mean. $N=20{,}000$.}\label{tab:shift}
\resizebox{\linewidth}{!}{\begin{tabular}{ll rrrrr}
\toprule
setting & $n$ & classical & PPI & PPI++ & lin & auto\\
\midrule
no shift (2013 only) & 100 & 92.3 & 95.2 & 91.7 & 95.6 & 95.2\\
 & 500 & 94.3 & 95.7 & 94.2 & 94.8 & 94.7\\
 & 2000 & 94.8 & 94.5 & 94.0 & 94.3 & 94.5\\
 & 10000 & 94.0 & 95.1 & 94.6 & 94.7 & 94.7\\
\midrule
shift 2013$\to$2014 & 100 & 53.7 & 94.2 & 94.1 & 96.5 & 96.3\\
 & 500 & 0.4 & 94.8 & 94.8 & 84.2 & 84.4\\
 & 2000 & 0.0 & 94.2 & 94.2 & 9.1 & 8.9\\
 & 10000 & 0.0 & 91.2 & 0.5 & 0.0 & 0.0\\
\bottomrule
\end{tabular}}
\end{table}

\subsection{Side information beyond the predictor}\label{sec:exp-x}
For the three census data sets the features $X$ are available. Using $X$ in the calibration (effectively a new predictor trained on the labels) is outside the fixed-predictor constraint, but the bound shows how much it could give. Cross-validated gradient-boosting fits of $Y$ on $(f,X)$ lower the residual variance by at most 0.9\% (Census income; 0.1\% and 0.3\% for the other two) relative to a fit on $f$ alone (Appendix Table~\ref{tab:xgain}), since $f$ already summarizes $X$ (age/sex, income, etc.); there is no headroom here either. We did not run interval-level experiments for $X$-augmented calibration; the statement is about the population risk only.

\subsection{Validity summary}
Across all real and synthetic cells with $n\ge100$ (plankton no-shift cells included, shift cells excluded) Table~\ref{tab:valid} shows that no method is systematically invalid; coverage is significantly below 95\% (Wilson upper limit $<95\%$) in 12 of 157 cells for classical and the cross-fitted lin, 14 for iso/bin10/auto, but 22--26 for PPI, PPI++ (both versions); the worst cells ($\approx$ 87--91\%) are Census income at $n=100$ (all methods). Because the replicates are shared across methods these counts are strongly correlated, and part of the flags are chance (about 4 of 157 expected under exact 95\% coverage), so they should be read as ``the cross-fitted intervals are at least as well calibrated as the PPI family'', not as exact validity. Cells above 95\% are rare (1--7).

\begin{table}[t]\centering\small
\caption{Coverage over all cells with $n\ge100$ (real, distorted, synthetic and plankton no-shift; shift cells excluded).}\label{tab:valid}
\resizebox{\linewidth}{!}{\begin{tabular}{lrrrrr}
\toprule
method & cells & mean cov. (\%) & min cov. (\%) & \#cells sig.\ below 95 & \#cells sig.\ above 95\\
\midrule
classical & 157 & 94.64 & 90.9 & 12 & 1\\
PPI & 157 & 94.56 & 88.5 & 22 & 0\\
PPI++ (ppi\_py) & 157 & 94.51 & 87.7 & 24 & 2\\
PPI++ (unclipped) & 157 & 94.46 & 87.1 & 26 & 2\\
cross-fit linear & 157 & 94.80 & 88.9 & 12 & 2\\
cross-fit isotonic & 157 & 94.78 & 89.0 & 14 & 7\\
cross-fit 10 bins & 157 & 94.81 & 88.6 & 14 & 3\\
cross-fit spline & 157 & 94.77 & 88.4 & 13 & 2\\
cross-fit auto & 157 & 94.80 & 89.1 & 14 & 3\\
\bottomrule
\end{tabular}}
\end{table}

\section{Discussion and limitations}\label{sec:disc}
\textbf{Answer to the question.} With labels, unlabeled data and predictor fixed, and the predictor output as the only information, an interval can be shorter than PPI++ only to the extent that $\mathbb E[Y\mid f]$ is not affine, by at most $1-\sqrt{1-sN/(n+N)}$ in width (Cor.~\ref{cor:gap}), and only if $n$ is large enough to learn the curve. On the seven data sets that ship with \texttt{ppi-python} the unmodified predictors are close to calibrated ($s\le2.5\%$), so the available gain is $\le1.2\%$ in width and cross-fitted recalibration did not realise it at $n\ge100$ (it cost 0.5--2\% at $n=500$ and up to 19\% at $n=50$). For predictors whose scores are monotone but not on the scale of the label the gain is real (2--27\% here) and is predicted accurately by the bound; there the improvement over PPI ($\lambda=1$) is much larger still. A hard 0/1 predictor leaves no room. Recalibration also removes the dependence on the scale of the scores, which PPI/PPI++ do not have. A practical rule that follows from the theory: plot the labeled $(f,Y)$ calibration curve, estimate $s$ and the bound by cross-validation, and recalibrate only if $(N/m)\,s$ is of the order of the estimation cost ($\approx$ a few percent at $n\approx500$); and use the unclipped PPI++ weight, which costs nothing and avoided losses of 2--6\% on two data sets.

\textbf{Limitations.}
(1) \emph{Target.} Only the population mean was studied; the closed-form gap uses the mean's influence function. For quantiles or regression coefficients the efficient estimators involve other conditional functions, and our numerical findings do not transfer without new checks.
(2) \emph{Information set.} The bound is for estimators that use only the scalar predictor output. Features $X$ could help in principle; we only measured population-level gains with gradient boosting on the three census sets (negligible), not interval-level behaviour, and learning from $X$ is closer to training a new predictor than to the fixed-predictor constraint of the task.
(3) \emph{Labeling mechanism.} Labeled points are an i.i.d.\ random subset (MCAR). Under shift (plankton) the recalibrated intervals were much less robust than PPI; we did not study covariate-shift corrections.
(4) \emph{Asymptotics and small $n$.} Intervals are Wald-type; Prop.~\ref{prop:valid} is asymptotic and not proved for the CV-selected learner (\textbf{auto}). Coverage at $n\le100$ is below nominal for all methods on heavy-tailed outcomes (Census income: 82--92\%) and can be 1--4 points low even on benign data; we do not claim to fix that, and ``shorter'' claims are only meaningful where coverage is nominal. A $t$-quantile did not repair PPI++ at $n=50$.
(5) \emph{Resampled populations.} Real-data cells resample the empirical distribution of finite example sets; for Forest (1596 records) and for $n$ comparable to $P$, repeated records make the effective problem easier for nonparametric fits. The distortions in Sec.~\ref{sec:exp-distort} are constructed, so we have no estimate of how often real deployed predictors are non-affinely calibrated; the two data sets with the largest $s$ (Gene expression 2.5\%, Census income 0.5\%) are those whose predictors are regressions on continuous outcomes, which is suggestive but not a sample of predictors.
(6) \emph{Novelty.} The efficiency bound and the cross-fitted augmented estimator are classical results in semiparametric/semi-supervised statistics; our contributions are the closed-form comparison with PPI and PPI++, the validity statement without rates, and the empirical map of when the gap is nonzero. Related prior work on the same estimand may exist that we did not find.
(7) \emph{Scope of evidence.} 2000 (synthetic) or 1000 (real) replicates per cell give coverage standard errors of about 0.5--0.7 points, so differences below one point in coverage are not resolved; width ratios are paired and much more precise.
\textbf{Broader impact.} The work is methodological, uses public data and has no foreseeable direct negative societal impact; the main risk is overconfidence from intervals used outside their assumptions (labeling mechanism, shift), which the failure cases document.


\bibliographystyle{plainnat}
\bibliography{refs}

\appendix
\section{Proofs}\label{app:proof}
\paragraph{Proposition~\ref{prop:ppi}.} Since $\mathbb E[Y-g(Z)\mid Z]=0$, $\mathrm{Var}(Y-\lambda Z)=\sigma^2+\mathrm{Var}(g(Z)-\lambda Z)$. Decompose $g(Z)=\mathbb E g+\beta(Z-\mathbb EZ)+h(Z)$ with $h\perp\{1,Z\}$ in $L^2$ (by the definition of $\beta$), so $\mathrm{Var}(g-\lambda Z)=v_h+(\beta-\lambda)^2V$. Hence $V_\lambda=\frac{\sigma^2+v_h}{n}+V\big[\frac{(\beta-\lambda)^2}{n}+\frac{\lambda^2}{N}\big]$. The bracket is a quadratic in $\lambda$ with leading coefficient $\frac1n+\frac1N=\frac m{nN}$, minimized at $\lambda^\star=\beta N/m$, where $\frac{(\beta-\lambda^\star)^2}{n}+\frac{\lambda^{\star2}}{N}=\beta^2\big(\frac{n}{m^2}+\frac{N}{m^2}\big)=\frac{\beta^2}{m}$. The formula for $V_\lambda-V_{\rm PPI++}$ is the quadratic's expansion around its minimizer.

\paragraph{Corollary~\ref{cor:gap}.} $V_{\rm PPI++}-V_{\rm eff}=\frac{v_h}{n}-\frac{v_h}{m}=v_h\frac{N}{nm}$. Dividing by $V_{\rm PPI++}=\frac{(\sigma^2+v_h)(1+\kappa)}{n}$ gives the ratio. For a two-point $Z$ every function of $Z$ is affine, so $h\equiv0$.

\paragraph{Proposition~\ref{prop:valid}.} Let $\tilde\theta=\frac1n\sum_L(Y_i-g^*(Z_i))+\frac1m\sum_{i\le m}g^*(Z_i)$ be the oracle estimator. Because $L$ is a uniformly random subset of an i.i.d.\ sample, $\tilde\theta-\theta$ is a sum of two uncorrelated terms (with $W_i=Y_i-g^*(Z_i)$, $\tilde\theta-\theta=\bar W_L-\mathbb EW+\ldots$); its variance is $\mathbb E(Y-g^*)^2/n+\mathrm{Var}\,g^*(Z)/m$ up to $O(1/m)$-relative terms and a Lindeberg-type CLT for sampling without replacement from an i.i.d.\ sample applies. Write $\delta_k=\hat g_{-k}-g^*$. Then
\[
\hat\theta-\tilde\theta=\sum_{k=1}^K\Big[\Big(\tfrac1m-\tfrac1n\Big)\sum_{i\in L_k}\delta_k(Z_i)+\tfrac1{Km}\sum_{j\notin L}\delta_k(Z_j)\Big]=:\sum_k D_k .
\]
Conditional on the training data for fold $k$, the points $Z_i,i\in L_k$ ($|L_k|=n/K$) and $Z_j,j\notin L$ ($N$ points) are i.i.d.\ and independent of $\delta_k$, so $\mathbb E[D_k\mid\text{train}]=\mathbb E\delta_k(Z)\big[(\tfrac1m-\tfrac1n)\tfrac nK+\tfrac N{Km}\big]=0$ (exact unbiasedness for any $\hat g$) and $\mathrm{Var}(D_k\mid\text{train})\le C\,\|\delta_k\|_{L^2}^2/n$ with $C$ depending only on $n/m$. By conditional Chebyshev, $D_k=o_p(n^{-1/2})$ when $\|\delta_k\|_{L^2}\to_p0$; no rate is needed. Hence $\hat\theta-\theta=(\tilde\theta-\theta)+o_p(n^{-1/2})$. The same argument applied to the sample variances (using $\widehat{\mathrm{Var}}_L(Y-\hat g)=\widehat{\mathrm{Var}}_L(Y-g^*)+o_p(1)$, which follows from $\|\delta_k\|_{L^2}\to_p 0$ and Cauchy--Schwarz, and similarly for $\widehat{\mathrm{Var}}_{\rm all}(\hat g)$) shows $\hat V/V^*\to_p1$. Slutsky's lemma finishes the proof. If $g^*=g$, $V^*=V_{\rm eff}$ by Proposition~\ref{prop:bound}. \qed

\paragraph{Remark (what is \emph{not} proved).} The proposition requires $\hat g_{-k}$ to converge to a deterministic limit. For the data-driven choice among lin/iso/spl this holds when the candidates' limits have different risks, but not necessarily when two of them are tied with different limits; we do not prove coverage of \textbf{auto} and evaluate it by simulation. All results are asymptotic in $n$ and $N$; no finite-sample guarantee is claimed.

\section{Additional results}\label{app:extra}
\begin{table}[h]\centering\small
\caption{Logit-scale synthetic classifier ($Y\sim\mathrm{Bernoulli}(\mathrm{expit}(f))$, $f\sim N(-1,\tau^2)$; top $\tau=2$, bottom $\tau=4$); 2000 replicates. Widths relative to PPI++ (unclipped); bound from the closed form with exact moments.}\label{tab:logit}
\resizebox{\linewidth}{!}{\begin{tabular}{lrr rrrr rrr}
\toprule
logit scores & $n$ & $N$ & \multicolumn{4}{c}{width / PPI++ (unclipped)} & \multicolumn{3}{c}{coverage (\%)}\\
\cmidrule(lr){4-7}\cmidrule(lr){8-10}
 & & & PPI & auto & iso & bound & PPI++ & iso & auto\\
\midrule
$s{=}0.06$ & 50 & 20000 & 4.58 & 1.055$\;{\scriptstyle\pm0.005}$ & 1.041$\;{\scriptstyle\pm0.003}$ & 0.970 & 93.3 & 93.8 & 94.2\\
 & 200 & 20000 & 4.58 & 0.996$\;{\scriptstyle\pm0.001}$ & 0.997$\;{\scriptstyle\pm0.001}$ & 0.971 & 94.8 & 94.7 & 94.5\\
 & 500 & 500 & 6.14 & 1.001$\;{\scriptstyle\pm0.000}$ & 1.011$\;{\scriptstyle\pm0.000}$ & 0.988 & 94.8 & 95.2 & 95.0\\
 & 500 & 5000 & 4.73 & 0.984$\;{\scriptstyle\pm0.000}$ & 0.990$\;{\scriptstyle\pm0.001}$ & 0.974 & 95.8 & 96.0 & 95.6\\
 & 500 & 20000 & 4.60 & 0.981$\;{\scriptstyle\pm0.000}$ & 0.986$\;{\scriptstyle\pm0.001}$ & 0.971 & 95.5 & 96.0 & 95.5\\
 & 500 & 50000 & 4.57 & 0.981$\;{\scriptstyle\pm0.000}$ & 0.985$\;{\scriptstyle\pm0.001}$ & 0.971 & 94.8 & 95.0 & 95.0\\
 & 3000 & 20000 & 4.82 & 0.978$\;{\scriptstyle\pm0.000}$ & 0.981$\;{\scriptstyle\pm0.000}$ & 0.976 & 95.6 & 95.5 & 95.6\\
\midrule
$s{=}0.23$ & 50 & 20000 & 10.65 & 0.967$\;{\scriptstyle\pm0.006}$ & 0.946$\;{\scriptstyle\pm0.005}$ & 0.877 & 94.0 & 94.2 & 94.0\\
 & 200 & 20000 & 10.74 & 0.904$\;{\scriptstyle\pm0.002}$ & 0.905$\;{\scriptstyle\pm0.002}$ & 0.879 & 95.3 & 95.2 & 95.2\\
 & 500 & 500 & 12.83 & 0.978$\;{\scriptstyle\pm0.000}$ & 0.983$\;{\scriptstyle\pm0.000}$ & 0.962 & 94.3 & 95.0 & 94.7\\
 & 500 & 5000 & 10.88 & 0.914$\;{\scriptstyle\pm0.001}$ & 0.914$\;{\scriptstyle\pm0.001}$ & 0.899 & 94.3 & 94.8 & 94.7\\
 & 500 & 20000 & 10.79 & 0.899$\;{\scriptstyle\pm0.001}$ & 0.898$\;{\scriptstyle\pm0.001}$ & 0.883 & 95.0 & 95.2 & 95.2\\
 & 500 & 50000 & 10.76 & 0.895$\;{\scriptstyle\pm0.001}$ & 0.894$\;{\scriptstyle\pm0.001}$ & 0.879 & 94.5 & 95.2 & 95.5\\
 & 3000 & 20000 & 10.93 & 0.913$\;{\scriptstyle\pm0.000}$ & 0.913$\;{\scriptstyle\pm0.000}$ & 0.908 & 94.9 & 95.1 & 95.0\\
\bottomrule
\end{tabular}}
\end{table}
\begin{table}[h]\centering\small
\caption{Ballots application (real fixed sets; no ground truth). Cross-fitted rows average the width over 20 fold-assignment seeds (s.d.\ over seeds $\le 2\times10^{-5}$).}\label{tab:ballots}
\resizebox{\linewidth}{!}{\begin{tabular}{lrrr}
\toprule
method & lower & upper & width\\
\midrule
classical & 0.5985 & 0.6576 & 0.05912\\
PPI & 0.6145 & 0.6278 & 0.01329\\
PPI++ (ppi\_py) & 0.6149 & 0.6280 & 0.01312\\
PPI++ (unclipped) & 0.6149 & 0.6280 & 0.01312\\
cross-fit linear & 0.6149 & 0.6280 & 0.01315\\
cross-fit isotonic & 0.6149 & 0.6280 & 0.01315\\
cross-fit 10 bins & 0.6149 & 0.6280 & 0.01315\\
cross-fit spline & 0.6149 & 0.6280 & 0.01315\\
cross-fit auto & 0.6149 & 0.6280 & 0.01315\\
\bottomrule
\end{tabular}}
\end{table}
\begin{table}[h]\centering\small
\caption{Residual variance of $Y$ (relative to $\mathrm{Var}(Y)$) from 5-fold cross-validated fits on the census sets: linear in $f$, gradient boosting (GB) on $f$, on $X$, on $(f,X)$. The last column is the relative reduction from adding $X$ to $f$.}\label{tab:xgain}
\resizebox{\linewidth}{!}{\begin{tabular}{lrrrrr}
\toprule
data set & lin($f$) & GB($f$) & GB($X$) & GB($f,X$) & gain of $X$ over GB($f$)\\
\midrule
Census income & 0.6309 & 0.6283 & 0.8474 & 0.6254 & 0.47\%\\
Census health ins. & 0.5004 & 0.5005 & 0.8428 & 0.5000 & 0.09\%\\
Census education & 0.6079 & 0.6080 & 0.8457 & 0.6062 & 0.29\%\\
\bottomrule
\end{tabular}}
\end{table}
\begin{table}[h]\centering\scriptsize
\caption{Real data, shipped predictors, $N=20{,}000$, varying $n$ (1000 replicates). Width ratios are relative to PPI++ (unclipped).}\label{tab:app_real}
\resizebox{\linewidth}{!}{\begin{tabular}{lrrrrrrr}
\toprule
 & $n{=}50$ & $n{=}100$ & $n{=}200$ & $n{=}500$ & $n{=}1000$ & $n{=}2000$ & $n{=}5000$\\
\midrule
\multicolumn{8}{l}{\emph{width of auto / width of PPI++ (unclipped)}}\\
Galaxies & 1.058 & 1.031 & 1.016 & 1.007 & 1.004 & 1.002 & 1.001\\
Forest & 1.146 & 1.032 & 1.014 & 1.005 & 0.999 & 0.994 & 0.991\\
AlphaFold & 1.027 & 1.014 & 1.011 & 1.005 & 1.003 & 1.002 & 1.000\\
Gene expression & 0.976 & 0.987 & 0.992 & 0.994 & 0.994 & 0.996 & 0.998\\
Census income & 1.191 & 1.085 & 1.043 & 1.021 & 1.011 & 1.006 & 1.002\\
Census health ins. & 1.064 & 1.035 & 1.018 & 1.008 & 1.004 & 1.002 & 1.001\\
Census education & 1.113 & 1.049 & 1.027 & 1.011 & 1.005 & 1.003 & 1.001\\
\midrule
\multicolumn{8}{l}{\emph{width of PPI ($\lambda{=}1$) / width of PPI++ (unclipped)}}\\
Galaxies & 0.984 & 0.992 & 0.996 & 1.000 & 1.002 & 1.007 & 1.030\\
Forest & 0.984 & 0.994 & 0.999 & 1.003 & 1.005 & 1.009 & 1.023\\
AlphaFold & 0.943 & 0.971 & 0.985 & 0.994 & 0.998 & 1.002 & 1.024\\
Gene expression & 1.017 & 1.055 & 1.060 & 1.049 & 1.027 & 1.004 & 1.010\\
Census income & 1.012 & 1.019 & 1.019 & 1.018 & 1.015 & 1.010 & 1.002\\
Census health ins. & 0.996 & 0.999 & 0.999 & 1.000 & 1.002 & 1.006 & 1.024\\
Census education & 1.000 & 1.000 & 1.001 & 1.001 & 1.002 & 1.004 & 1.018\\
\midrule
\multicolumn{8}{l}{\emph{coverage (\%) of auto}}\\
Galaxies & 94.2 & 95.3 & 95.4 & 95.2 & 95.1 & 95.3 & 94.9\\
Forest & 92.3 & 94.0 & 94.5 & 95.1 & 94.6 & 95.2 & 95.9\\
AlphaFold & 89.2 & 93.3 & 94.3 & 94.2 & 94.3 & 95.1 & 94.2\\
Gene expression & 95.6 & 94.7 & 94.3 & 93.8 & 95.2 & 95.3 & 94.0\\
Census income & 84.8 & 89.3 & 91.7 & 93.4 & 94.2 & 94.8 & 95.3\\
Census health ins. & 94.7 & 94.5 & 95.8 & 95.0 & 95.0 & 96.3 & 94.2\\
Census education & 93.6 & 93.1 & 93.3 & 94.5 & 94.8 & 94.5 & 94.7\\
\midrule
\multicolumn{8}{l}{\emph{coverage (\%) of PPI++ (unclipped)}}\\
Galaxies & 92.8 & 93.7 & 95.1 & 94.7 & 94.7 & 95.3 & 94.7\\
Forest & 89.5 & 92.9 & 93.5 & 95.1 & 94.9 & 95.1 & 95.9\\
AlphaFold & 86.6 & 92.3 & 91.9 & 94.3 & 94.0 & 95.0 & 94.1\\
Gene expression & 93.2 & 93.2 & 93.8 & 94.0 & 95.2 & 95.2 & 94.5\\
Census income & 81.6 & 87.1 & 90.4 & 93.1 & 93.4 & 94.7 & 95.5\\
Census health ins. & 94.3 & 93.9 & 95.6 & 94.8 & 95.1 & 96.1 & 94.3\\
Census education & 91.7 & 92.4 & 93.0 & 94.1 & 94.8 & 94.7 & 94.6\\
\bottomrule
\end{tabular}}
\end{table}
\begin{table}[h]\centering\scriptsize
\caption{Real data, shipped predictors, $n=500$, varying the number of unlabeled points $N$ (1000 replicates).}\label{tab:app_N}
\resizebox{\linewidth}{!}{\begin{tabular}{lrrrrr}
\toprule
 & $N{=}50$ & $N{=}250$ & $N{=}1000$ & $N{=}5000$ & $N{=}50000$\\
\midrule
\multicolumn{6}{l}{\emph{width of auto / width of PPI++ (unclipped), $n=500$}}\\
Galaxies & 1.009 & 1.009 & 1.008 & 1.007 & 1.007\\
Forest & 1.019 & 1.016 & 1.011 & 1.006 & 1.004\\
AlphaFold & 1.011 & 1.011 & 1.010 & 1.006 & 1.003\\
Gene expression & 1.003 & 1.003 & 1.002 & 0.998 & 0.991\\
Census income & 1.025 & 1.024 & 1.021 & 1.021 & 1.021\\
Census health ins. & 1.009 & 1.009 & 1.008 & 1.008 & 1.007\\
Census education & 1.014 & 1.014 & 1.012 & 1.011 & 1.011\\
\midrule
\multicolumn{6}{l}{\emph{width of PPI ($\lambda{=}1$) / width of PPI++ (unclipped)}}\\
Galaxies & 2.507 & 1.402 & 1.079 & 1.007 & 0.999\\
Forest & 2.060 & 1.264 & 1.053 & 1.008 & 1.001\\
AlphaFold & 2.582 & 1.423 & 1.079 & 0.999 & 0.994\\
Gene expression & 2.753 & 1.468 & 1.072 & 1.002 & 1.073\\
Census income & 1.662 & 1.102 & 1.005 & 1.010 & 1.020\\
Census health ins. & 2.421 & 1.355 & 1.066 & 1.005 & 1.000\\
Census education & 2.200 & 1.287 & 1.052 & 1.005 & 1.000\\
\midrule
\multicolumn{6}{l}{\emph{width of classical / width of PPI++ (unclipped)}}\\
Galaxies & 1.027 & 1.106 & 1.259 & 1.416 & 1.482\\
Forest & 1.014 & 1.053 & 1.115 & 1.165 & 1.186\\
AlphaFold & 1.033 & 1.135 & 1.338 & 1.590 & 1.703\\
Gene expression & 1.046 & 1.204 & 1.619 & 2.519 & 3.480\\
Census income & 1.018 & 1.070 & 1.158 & 1.238 & 1.269\\
Census health ins. & 1.024 & 1.096 & 1.225 & 1.357 & 1.407\\
Census education & 1.019 & 1.073 & 1.165 & 1.249 & 1.279\\
\midrule
\multicolumn{6}{l}{\emph{coverage (\%) of auto}}\\
Galaxies & 94.9 & 94.8 & 96.0 & 96.0 & 95.2\\
Forest & 96.8 & 94.6 & 95.7 & 93.8 & 94.3\\
AlphaFold & 95.1 & 95.6 & 94.9 & 95.7 & 95.6\\
Gene expression & 95.9 & 95.7 & 96.4 & 95.7 & 95.3\\
Census income & 95.8 & 94.2 & 93.6 & 92.8 & 94.1\\
Census health ins. & 95.6 & 95.1 & 95.4 & 94.6 & 94.4\\
Census education & 93.9 & 93.6 & 94.6 & 95.7 & 94.9\\
\bottomrule
\end{tabular}}
\end{table}
\begin{figure}[h]\centering
\includegraphics[width=\linewidth]{../figs/fig_real.pdf}
\caption{Real data, shipped predictors: width of cross-fitted intervals relative to PPI++ (unclipped) versus $n$ ($N=20{,}000$; $\pm1.96$ s.e.). Dashed: bound estimated from the population. Classical and PPI curves are off this scale (Table~\ref{tab:app_real}). The Forest population has 1596 records, so $n\ge1000$ is affected by repeated records.}
\end{figure}
\begin{figure}[h]\centering
\includegraphics[width=\linewidth]{../figs/fig_distort.pdf}
\caption{Distorted predictor outputs: same quantities as above ($N=20{,}000$).}
\end{figure}

\section{Reproducibility details}\label{app:repro}
\textbf{Software and data.} Python 3.11.2, numpy 2.4.6, scipy 1.17.1, pandas 3.1.1, scikit-learn 1.9.1, matplotlib 3.11.2, \texttt{ppi-python} 0.2.3 (MIT; the \texttt{.npz} example data carry no separate license statement in the repository; original sources are cited in its example notebooks), gdown 6.5.1. The data are downloaded by \texttt{ppi\_py.datasets.load\_dataset}. \textbf{Seeds.} Replicate $r$ of cell $c$ uses \texttt{SeedSequence([20261111, crc32(c), r])}. \textbf{Compute.} CPU only; the container exposes 4 effective cores (cgroup quota) on a shared host; the full set of experiments (E1a--E4) took about 2.5 wall-clock hours, i.e.\ roughly 10 CPU-hours. Per-experiment commands and the mapping from every table and figure to its script are in the supplementary \texttt{README.md}. Preliminary runs (a 40-replicate smoke test of the runner and an aborted first launch that was restarted after replacing a slow spline basis by a faster natural-spline implementation) are not reported; no other experiments were run and discarded.

\newpage
\section*{NeurIPS Paper Checklist}

%%% BEGIN INSTRUCTIONS %%%
%%% END INSTRUCTIONS %%%


\begin{enumerate}

\item {\bf Claims}
    \item[] Question: Do the main claims made in the abstract and introduction accurately reflect the paper's contributions and scope?
    \item[] Answer: \answerYes{} % Replace by \answerYes{}, \answerNo{}, or \answerNA{}.
    \item[] Justification: Abstract and introduction state the theory (bound, Cor.~\ref{cor:gap}), the estimator, and the empirical findings including the negative result on the unmodified predictors; scope limited to the mean with a scalar predictor output.
    \item[] Guidelines:
    \begin{itemize}
        \item The answer \answerNA{} means that the abstract and introduction do not include the claims made in the paper.
        \item The abstract and/or introduction should clearly state the claims made, including the contributions made in the paper and important assumptions and limitations. A \answerNo{} or \answerNA{} answer to this question will not be perceived well by the reviewers. 
        \item The claims made should match theoretical and experimental results, and reflect how much the results can be expected to generalize to other settings. 
        \item It is fine to include aspirational goals as motivation as long as it is clear that these goals are not attained by the paper. 
    \end{itemize}

\item {\bf Limitations}
    \item[] Question: Does the paper discuss the limitations of the work performed by the authors?
    \item[] Answer: \answerYes{} % Replace by \answerYes{}, \answerNo{}, or \answerNA{}.
    \item[] Justification: Section~\ref{sec:disc} (Limitations) lists assumptions (MCAR labeling, scalar predictor output, asymptotic Wald intervals, mean target only, resampled finite populations, synthetic distortions) and failure cases (small $n$, binary predictor, distribution shift).
    \item[] Guidelines:
    \begin{itemize}
        \item The answer \answerNA{} means that the paper has no limitation while the answer \answerNo{} means that the paper has limitations, but those are not discussed in the paper. 
        \item The authors are encouraged to create a separate ``Limitations'' section in their paper.
        \item The paper should point out any strong assumptions and how robust the results are to violations of these assumptions (e.g., independence assumptions, noiseless settings, model well-specification, asymptotic approximations only holding locally). The authors should reflect on how these assumptions might be violated in practice and what the implications would be.
        \item The authors should reflect on the scope of the claims made, e.g., if the approach was only tested on a few datasets or with a few runs. In general, empirical results often depend on implicit assumptions, which should be articulated.
        \item The authors should reflect on the factors that influence the performance of the approach. For example, a facial recognition algorithm may perform poorly when image resolution is low or images are taken in low lighting. Or a speech-to-text system might not be used reliably to provide closed captions for online lectures because it fails to handle technical jargon.
        \item The authors should discuss the computational efficiency of the proposed algorithms and how they scale with dataset size.
        \item If applicable, the authors should discuss possible limitations of their approach to address problems of privacy and fairness.
        \item While the authors might fear that complete honesty about limitations might be used by reviewers as grounds for rejection, a worse outcome might be that reviewers discover limitations that aren't acknowledged in the paper. The authors should use their best judgment and recognize that individual actions in favor of transparency play an important role in developing norms that preserve the integrity of the community. Reviewers will be specifically instructed to not penalize honesty concerning limitations.
    \end{itemize}

\item {\bf Theory assumptions and proofs}
    \item[] Question: For each theoretical result, does the paper provide the full set of assumptions and a complete (and correct) proof?
    \item[] Answer: \answerYes{} % Replace by \answerYes{}, \answerNo{}, or \answerNA{}.
    \item[] Justification: Assumptions are in Sec.~\ref{sec:theory}--\ref{sec:method}; proofs (or standard-result sketches with references) are in Appendix~\ref{app:proof}. Coverage of the CV-selected learner is explicitly \emph{not} proved.
    \item[] Guidelines:
    \begin{itemize}
        \item The answer \answerNA{} means that the paper does not include theoretical results. 
        \item All the theorems, formulas, and proofs in the paper should be numbered and cross-referenced.
        \item All assumptions should be clearly stated or referenced in the statement of any theorems.
        \item The proofs can either appear in the main paper or the supplemental material, but if they appear in the supplemental material, the authors are encouraged to provide a short proof sketch to provide intuition. 
        \item Inversely, any informal proof provided in the core of the paper should be complemented by formal proofs provided in appendix or supplemental material.
        \item Theorems and Lemmas that the proof relies upon should be properly referenced. 
    \end{itemize}

    \item {\bf Experimental result reproducibility}
    \item[] Question: Does the paper fully disclose all the information needed to reproduce the main experimental results of the paper to the extent that it affects the main claims and/or conclusions of the paper (regardless of whether the code and data are provided or not)?
    \item[] Answer: \answerYes{} % Replace by \answerYes{}, \answerNo{}, or \answerNA{}.
    \item[] Justification: Sec.~\ref{sec:exp} and Appendix~\ref{app:repro} give data, resampling protocol, sample sizes, learners and hyperparameters, seeds and the number of replicates; all code is supplied.
    \item[] Guidelines:
    \begin{itemize}
        \item The answer \answerNA{} means that the paper does not include experiments.
        \item If the paper includes experiments, a \answerNo{} answer to this question will not be perceived well by the reviewers: Making the paper reproducible is important, regardless of whether the code and data are provided or not.
        \item If the contribution is a dataset and\slash or model, the authors should describe the steps taken to make their results reproducible or verifiable. 
        \item Depending on the contribution, reproducibility can be accomplished in various ways. For example, if the contribution is a novel architecture, describing the architecture fully might suffice, or if the contribution is a specific model and empirical evaluation, it may be necessary to either make it possible for others to replicate the model with the same dataset, or provide access to the model. In general. releasing code and data is often one good way to accomplish this, but reproducibility can also be provided via detailed instructions for how to replicate the results, access to a hosted model (e.g., in the case of a large language model), releasing of a model checkpoint, or other means that are appropriate to the research performed.
        \item While NeurIPS does not require releasing code, the conference does require all submissions to provide some reasonable avenue for reproducibility, which may depend on the nature of the contribution. For example
        \begin{enumerate}
            \item If the contribution is primarily a new algorithm, the paper should make it clear how to reproduce that algorithm.
            \item If the contribution is primarily a new model architecture, the paper should describe the architecture clearly and fully.
            \item If the contribution is a new model (e.g., a large language model), then there should either be a way to access this model for reproducing the results or a way to reproduce the model (e.g., with an open-source dataset or instructions for how to construct the dataset).
            \item We recognize that reproducibility may be tricky in some cases, in which case authors are welcome to describe the particular way they provide for reproducibility. In the case of closed-source models, it may be that access to the model is limited in some way (e.g., to registered users), but it should be possible for other researchers to have some path to reproducing or verifying the results.
        \end{enumerate}
    \end{itemize}


\item {\bf Open access to data and code}
    \item[] Question: Does the paper provide open access to the data and code, with sufficient instructions to faithfully reproduce the main experimental results, as described in supplemental material?
    \item[] Answer: \answerYes{} % Replace by \answerYes{}, \answerNo{}, or \answerNA{}.
    \item[] Justification: All data are public (\texttt{ppi-python} v0.2.3 example sets) or synthetic; code, exact commands, versions and the table/figure-to-script mapping are in the supplementary README (anonymized).
    \item[] Guidelines:
    \begin{itemize}
        \item The answer \answerNA{} means that paper does not include experiments requiring code.
        \item Please see the NeurIPS code and data submission guidelines (\url{https://neurips.cc/public/guides/CodeSubmissionPolicy}) for more details.
        \item While we encourage the release of code and data, we understand that this might not be possible, so \answerNo{} is an acceptable answer. Papers cannot be rejected simply for not including code, unless this is central to the contribution (e.g., for a new open-source benchmark).
        \item The instructions should contain the exact command and environment needed to run to reproduce the results. See the NeurIPS code and data submission guidelines (\url{https://neurips.cc/public/guides/CodeSubmissionPolicy}) for more details.
        \item The authors should provide instructions on data access and preparation, including how to access the raw data, preprocessed data, intermediate data, and generated data, etc.
        \item The authors should provide scripts to reproduce all experimental results for the new proposed method and baselines. If only a subset of experiments are reproducible, they should state which ones are omitted from the script and why.
        \item At submission time, to preserve anonymity, the authors should release anonymized versions (if applicable).
        \item Providing as much information as possible in supplemental material (appended to the paper) is recommended, but including URLs to data and code is permitted.
    \end{itemize}


\item {\bf Experimental setting/details}
    \item[] Question: Does the paper specify all the training and test details (e.g., data splits, hyperparameters, how they were chosen, type of optimizer) necessary to understand the results?
    \item[] Answer: \answerYes{} % Replace by \answerYes{}, \answerNo{}, or \answerNA{}.
    \item[] Justification: Sec.~\ref{sec:exp} and Appendix~\ref{app:repro}: K=5, knots, bins, nested CV for the selector, sample sizes; nothing was tuned on evaluation replicates. No training of predictors is involved.
    \item[] Guidelines:
    \begin{itemize}
        \item The answer \answerNA{} means that the paper does not include experiments.
        \item The experimental setting should be presented in the core of the paper to a level of detail that is necessary to appreciate the results and make sense of them.
        \item The full details can be provided either with the code, in appendix, or as supplemental material.
    \end{itemize}

\item {\bf Experiment statistical significance}
    \item[] Question: Does the paper report error bars suitably and correctly defined or other appropriate information about the statistical significance of the experiments?
    \item[] Answer: \answerYes{} % Replace by \answerYes{}, \answerNo{}, or \answerNA{}.
    \item[] Justification: Coverage is reported with Wilson 95\% intervals, widths as paired ratios with $\pm1.96$ standard errors over replicates (2000 synthetic / 1000 real-data replicates per cell); the source of variability (resampling of labeled and unlabeled sets and fold assignment) is stated.
    \item[] Guidelines:
    \begin{itemize}
        \item The answer \answerNA{} means that the paper does not include experiments.
        \item The authors should answer \answerYes{} if the results are accompanied by error bars, confidence intervals, or statistical significance tests, at least for the experiments that support the main claims of the paper.
        \item The factors of variability that the error bars are capturing should be clearly stated (for example, train/test split, initialization, random drawing of some parameter, or overall run with given experimental conditions).
        \item The method for calculating the error bars should be explained (closed form formula, call to a library function, bootstrap, etc.)
        \item The assumptions made should be given (e.g., Normally distributed errors).
        \item It should be clear whether the error bar is the standard deviation or the standard error of the mean.
        \item It is OK to report 1-sigma error bars, but one should state it. The authors should preferably report a 2-sigma error bar than state that they have a 96\% CI, if the hypothesis of Normality of errors is not verified.
        \item For asymmetric distributions, the authors should be careful not to show in tables or figures symmetric error bars that would yield results that are out of range (e.g., negative error rates).
        \item If error bars are reported in tables or plots, the authors should explain in the text how they were calculated and reference the corresponding figures or tables in the text.
    \end{itemize}

\item {\bf Experiments compute resources}
    \item[] Question: For each experiment, does the paper provide sufficient information on the computer resources (type of compute workers, memory, time of execution) needed to reproduce the experiments?
    \item[] Answer: \answerYes{} % Replace by \answerYes{}, \answerNo{}, or \answerNA{}.
    \item[] Justification: CPU only; Appendix~\ref{app:repro} reports the machine (4 effective cores, shared host) and approximate wall-clock/CPU time per experiment.
    \item[] Guidelines:
    \begin{itemize}
        \item The answer \answerNA{} means that the paper does not include experiments.
        \item The paper should indicate the type of compute workers CPU or GPU, internal cluster, or cloud provider, including relevant memory and storage.
        \item The paper should provide the amount of compute required for each of the individual experimental runs as well as estimate the total compute. 
        \item The paper should disclose whether the full research project required more compute than the experiments reported in the paper (e.g., preliminary or failed experiments that didn't make it into the paper). 
    \end{itemize}
    
\item {\bf Code of ethics}
    \item[] Question: Does the research conducted in the paper conform, in every respect, with the NeurIPS Code of Ethics \url{https://neurips.cc/public/EthicsGuidelines}?
    \item[] Answer: \answerYes{} % Replace by \answerYes{}, \answerNo{}, or \answerNA{}.
    \item[] Justification: The research uses public data and simulation only and conforms to the NeurIPS Code of Ethics.
    \item[] Guidelines:
    \begin{itemize}
        \item The answer \answerNA{} means that the authors have not reviewed the NeurIPS Code of Ethics.
        \item If the authors answer \answerNo, they should explain the special circumstances that require a deviation from the Code of Ethics.
        \item The authors should make sure to preserve anonymity (e.g., if there is a special consideration due to laws or regulations in their jurisdiction).
    \end{itemize}


\item {\bf Broader impacts}
    \item[] Question: Does the paper discuss both potential positive societal impacts and negative societal impacts of the work performed?
    \item[] Answer: \answerNA{} % Replace by \answerYes{}, \answerNo{}, or \answerNA{}.
    \item[] Justification: This is foundational statistical methodology evaluated on public data sets; we do not foresee a direct societal impact beyond the generic risk that overly optimistic intervals could mislead downstream analyses, which is why we report failure cases (e.g.\ distribution shift).
    \item[] Guidelines:
    \begin{itemize}
        \item The answer \answerNA{} means that there is no societal impact of the work performed.
        \item If the authors answer \answerNA{} or \answerNo, they should explain why their work has no societal impact or why the paper does not address societal impact.
        \item Examples of negative societal impacts include potential malicious or unintended uses (e.g., disinformation, generating fake profiles, surveillance), fairness considerations (e.g., deployment of technologies that could make decisions that unfairly impact specific groups), privacy considerations, and security considerations.
        \item The conference expects that many papers will be foundational research and not tied to particular applications, let alone deployments. However, if there is a direct path to any negative applications, the authors should point it out. For example, it is legitimate to point out that an improvement in the quality of generative models could be used to generate Deepfakes for disinformation. On the other hand, it is not needed to point out that a generic algorithm for optimizing neural networks could enable people to train models that generate Deepfakes faster.
        \item The authors should consider possible harms that could arise when the technology is being used as intended and functioning correctly, harms that could arise when the technology is being used as intended but gives incorrect results, and harms following from (intentional or unintentional) misuse of the technology.
        \item If there are negative societal impacts, the authors could also discuss possible mitigation strategies (e.g., gated release of models, providing defenses in addition to attacks, mechanisms for monitoring misuse, mechanisms to monitor how a system learns from feedback over time, improving the efficiency and accessibility of ML).
    \end{itemize}
    
\item {\bf Safeguards}
    \item[] Question: Does the paper describe safeguards that have been put in place for responsible release of data or models that have a high risk for misuse (e.g., pre-trained language models, image generators, or scraped datasets)?
    \item[] Answer: \answerNA{} % Replace by \answerYes{}, \answerNo{}, or \answerNA{}.
    \item[] Justification: No models or scraped data with high misuse risk are released.
    \item[] Guidelines:
    \begin{itemize}
        \item The answer \answerNA{} means that the paper poses no such risks.
        \item Released models that have a high risk for misuse or dual-use should be released with necessary safeguards to allow for controlled use of the model, for example by requiring that users adhere to usage guidelines or restrictions to access the model or implementing safety filters. 
        \item Datasets that have been scraped from the Internet could pose safety risks. The authors should describe how they avoided releasing unsafe images.
        \item We recognize that providing effective safeguards is challenging, and many papers do not require this, but we encourage authors to take this into account and make a best faith effort.
    \end{itemize}

\item {\bf Licenses for existing assets}
    \item[] Question: Are the creators or original owners of assets (e.g., code, data, models), used in the paper, properly credited and are the license and terms of use explicitly mentioned and properly respected?
    \item[] Answer: \answerYes{} % Replace by \answerYes{}, \answerNo{}, or \answerNA{}.
    \item[] Justification: \texttt{ppi-python} v0.2.3 (MIT) is credited; its example data files carry no separate license statement in the repository, which we state in Appendix~\ref{app:repro}; the underlying sources (Galaxy Zoo 2, ACS PUMS, etc.) are cited where known.
    \item[] Guidelines:
    \begin{itemize}
        \item The answer \answerNA{} means that the paper does not use existing assets.
        \item The authors should cite the original paper that produced the code package or dataset.
        \item The authors should state which version of the asset is used and, if possible, include a URL.
        \item The name of the license (e.g., CC-BY 4.0) should be included for each asset.
        \item For scraped data from a particular source (e.g., website), the copyright and terms of service of that source should be provided.
        \item If assets are released, the license, copyright information, and terms of use in the package should be provided. For popular datasets, \url{paperswithcode.com/datasets} has curated licenses for some datasets. Their licensing guide can help determine the license of a dataset.
        \item For existing datasets that are re-packaged, both the original license and the license of the derived asset (if it has changed) should be provided.
        \item If this information is not available online, the authors are encouraged to reach out to the asset's creators.
    \end{itemize}

\item {\bf New assets}
    \item[] Question: Are new assets introduced in the paper well documented and is the documentation provided alongside the assets?
    \item[] Answer: \answerYes{} % Replace by \answerYes{}, \answerNo{}, or \answerNA{}.
    \item[] Justification: The code released with the paper is documented in a README (commands, environment, expected run times, outputs).
    \item[] Guidelines:
    \begin{itemize}
        \item The answer \answerNA{} means that the paper does not release new assets.
        \item Researchers should communicate the details of the dataset\slash code\slash model as part of their submissions via structured templates. This includes details about training, license, limitations, etc. 
        \item The paper should discuss whether and how consent was obtained from people whose asset is used.
        \item At submission time, remember to anonymize your assets (if applicable). You can either create an anonymized URL or include an anonymized zip file.
    \end{itemize}

\item {\bf Crowdsourcing and research with human subjects}
    \item[] Question: For crowdsourcing experiments and research with human subjects, does the paper include the full text of instructions given to participants and screenshots, if applicable, as well as details about compensation (if any)? 
    \item[] Answer: \answerNA{} % Replace by \answerYes{}, \answerNo{}, or \answerNA{}.
    \item[] Justification: No crowdsourcing or human-subject experiments were run by us (the example data contain labels produced by others and are used as-is).
    \item[] Guidelines:
    \begin{itemize}
        \item The answer \answerNA{} means that the paper does not involve crowdsourcing nor research with human subjects.
        \item Including this information in the supplemental material is fine, but if the main contribution of the paper involves human subjects, then as much detail as possible should be included in the main paper. 
        \item According to the NeurIPS Code of Ethics, workers involved in data collection, curation, or other labor should be paid at least the minimum wage in the country of the data collector. 
    \end{itemize}

\item {\bf Institutional review board (IRB) approvals or equivalent for research with human subjects}
    \item[] Question: Does the paper describe potential risks incurred by study participants, whether such risks were disclosed to the subjects, and whether Institutional Review Board (IRB) approvals (or an equivalent approval/review based on the requirements of your country or institution) were obtained?
    \item[] Answer: \answerNA{} % Replace by \answerYes{}, \answerNo{}, or \answerNA{}.
    \item[] Justification: No human-subject research was conducted.
    \item[] Guidelines:
    \begin{itemize}
        \item The answer \answerNA{} means that the paper does not involve crowdsourcing nor research with human subjects.
        \item Depending on the country in which research is conducted, IRB approval (or equivalent) may be required for any human subjects research. If you obtained IRB approval, you should clearly state this in the paper. 
        \item We recognize that the procedures for this may vary significantly between institutions and locations, and we expect authors to adhere to the NeurIPS Code of Ethics and the guidelines for their institution. 
        \item For initial submissions, do not include any information that would break anonymity (if applicable), such as the institution conducting the review.
    \end{itemize}

\item {\bf Declaration of LLM usage}
    \item[] Question: Does the paper describe the usage of LLMs if it is an important, original, or non-standard component of the core methods in this research? Note that if the LLM is used only for writing, editing, or formatting purposes and does \emph{not} impact the core methodology, scientific rigor, or originality of the research, declaration is not required.
    %this research? 
    \item[] Answer: \answerNA{} % Replace by \answerYes{}, \answerNo{}, or \answerNA{}.
    \item[] Justification: The research itself (literature search, analysis, code, experiments and writing) was carried out by an autonomous AI research agent (Anthropic Claude); this is disclosed in the supplementary README. LLMs are not a component of the proposed method.
    \item[] Guidelines:
    \begin{itemize}
        \item The answer \answerYes{} means that the core method development in this research does not involve LLMs as any important, original, or non-standard components.
        \item Please refer to our LLM policy in the NeurIPS handbook for what should or should not be described.
    \end{itemize}

\end{enumerate}
\end{document}
